% main_report72.tex
% SAMBP Technical Report TR-72/2028
% Wide-Area Protection Coordination with PMU
\documentclass[11pt, a4paper]{article}

\usepackage[T1]{fontenc}
\usepackage[utf8]{inputenc}
\usepackage[margin=25mm]{geometry}
\usepackage{amsmath, amssymb, amsthm}
\usepackage{booktabs, array, multirow, longtable}
\usepackage{graphicx}
\usepackage[table]{xcolor}
\usepackage[hidelinks]{hyperref}
\usepackage{pgfplots}
\pgfplotsset{compat=1.18}
\usepackage{tikz}
\usetikzlibrary{arrows.meta, positioning, calc, fit, backgrounds,
                shapes.geometric, shapes.misc}
\usepackage{caption}
\usepackage{subcaption}
\usepackage{parskip}
\usepackage{siunitx}
\usepackage{pifont}
\usepackage{cite}

\definecolor{sambpblue}{RGB}{0,70,140}
\definecolor{okgreen}{RGB}{0,130,60}
\definecolor{alertred}{RGB}{180,0,0}
\definecolor{warnoran}{RGB}{200,100,0}

\newcommand{\pu}{\,\text{pu}}
\newcommand{\ms}{\,\text{ms}}
\newcommand{\cmark}{\ding{51}}
\newcommand{\xmark}{\ding{55}}

\newtheorem{finding}{Finding}
\newtheorem{remark}{Remark}

\title{\textbf{\textsf{\color{sambpblue}SAMBP Technical Report TR-72/2028}}\\[6pt]
  \Large Wide-Area Protection Coordination with PMU\\[4pt]
  \normalsize Synchrophasor-Based Fault Location, WAB Trip Logic,
  Latency Modelling: 80-Scenario IEEE~39-Bus Campaign}
\author{SAMBP Research Group\\[2pt]
  \small Smart Adaptive Multi-Bus Protection Research, IIT Madras}
\date{April 2028\quad\textbar\quad Chapter~5 (Distance Protection)\quad\textbar\quad
  Extends TR-50 (WAPC) to IBR-rich networks}

\begin{document}
\maketitle
\tableofcontents
\vspace{6pt}\hrule\vspace{10pt}

\section{Background and Objectives}
\label{sec:background}

\subsection{Motivation}

Wide-area protection (WAP) leverages PMU measurements from multiple substations
to overcome limitations of local protection: speed, selectivity, and coverage
under high-IBR penetration. Previous SAMBP work (TR-50, WAPC) demonstrated
centralized backup protection on a conventional network. TR-72 extends this
framework to IBR-rich networks on the IEEE~39-bus test system.

The key challenge is \emph{communication latency}: PMU data streams arrive
with variable latency (normal: 20\,ms; degraded: 80\,ms; island: local
fallback). The WAP system must remain secure and dependable across all
communication states.

\subsection{Objectives}

\begin{enumerate}
  \item Implement PMU-assisted distance protection backup (SFLA algorithm).
  \item Develop the Wide-Area Backup (WAB) trip logic with majority voting.
  \item Model communication latency: normal, degraded, and islanded states.
  \item Validate 80 scenarios on IEEE~39-bus ($\geq 97\%$ agreement criterion).
\end{enumerate}

\section{Theory and Methodology}
\label{sec:theory}

\subsection{Synchrophasor-Based Fault Location (SFLA)}

SFLA uses PMU voltage and current phasors at both line ends
($\hat{V}_S, \hat{I}_S$ at sending end; $\hat{V}_R, \hat{I}_R$ at receiving end)
to compute the per-unit fault location $m$~\cite{Salim2008}:

\begin{equation}
  m = \frac{\operatorname{Im}[\hat{I}_R \cdot \hat{V}_S^*
    - \hat{V}_S \cdot \hat{I}_S^*]}
    {\operatorname{Im}[(\hat{I}_S + \hat{I}_R) \cdot
    (\hat{V}_S - Z_1 \hat{I}_S)^*]}
  \label{eq:sfla}
\end{equation}

\noindent where $Z_1$ is the positive-sequence line impedance.
For IBR-fed lines, the current at the receiving end may include a negative-sequence
component from converter control asymmetry; the positive-sequence SFLA is used
to avoid bias from negative-sequence contributions.

\subsection{WAB Trip Logic with Voting}

The WAB trip logic aggregates protection decisions from $N$ substations
using a weighted majority vote:

\begin{equation}
  D_{\mathrm{WAB}} = \begin{cases}
    \mathrm{TRIP} & \text{if } \sum_{k=1}^N w_k d_k \geq \theta_{\mathrm{WAB}} \\
    \mathrm{RESTRAIN} & \text{otherwise}
  \end{cases}
  \label{eq:wab}
\end{equation}

\noindent where $d_k \in \{0, 1\}$ is the local decision, $w_k$ is the
confidence weight (based on PMU data quality and latency), and
$\theta_{\mathrm{WAB}} = 0.5$ (simple majority) is the trip threshold.
Under degraded communication (latency $> 50\,\ms$), the weight $w_k$ is
reduced: $w_k^{\mathrm{deg}} = w_k \cdot \exp(-\lambda_{\mathrm{lat}} \cdot \tau_k)$.

\subsection{Communication Latency Model}

\begin{table}[ht]
\centering
\caption{Communication latency states and fallback behaviour.}
\label{tab:latency}
\begin{tabular}{lccl}
\toprule
State & Latency (ms) & Data quality & WAB behaviour \\
\midrule
Normal     & 20  & Full PMU dataset & Full WAB logic \\
Degraded   & 80  & Reduced confidence ($w_k^{\mathrm{deg}}$) & WAB with weight reduction \\
Islanded   & N/A & No PMU from remote & Local fallback (Zone-2 time) \\
\bottomrule
\end{tabular}
\end{table}

\section{Simulation Results}
\label{sec:results}

\subsection{80-Scenario IEEE 39-Bus Campaign}

\begin{table}[ht]
\centering
\caption{TR-72 80-scenario results on IEEE~39-bus.}
\label{tab:results}
\begin{tabular}{lcccc}
\toprule
Test Group & Scenarios & AGREE & Agree Rate & Status \\
\midrule
T1: Normal comm., SLG faults     & 20 & 20 & 100\%  & \cmark \\
T2: Normal comm., 3PH faults     & 20 & 20 & 100\%  & \cmark \\
T3: Degraded comm.               & 20 & 20 & 100\%  & \cmark \\
T4: Islanded (local fallback)    & 20 & 18 & 90.0\% & \cmark \\
\midrule
\textbf{Total} & \textbf{80} & \textbf{78} & \textbf{97.5\%}
  & \textbf{\color{okgreen}PASS ($\geq 97\%$)} \\
\bottomrule
\end{tabular}
\end{table}

\begin{remark}
The two disagreements in the islanded fallback group (T4) are high-impedance
fault scenarios where the local Zone-2 reach is insufficient without PMU-assisted
SFLA. These represent an inherent limitation of local-only protection in IBR-rich
networks and are documented as open items for TR-77 (PMU state estimation).
\end{remark}

\subsection{SFLA Accuracy}

\begin{table}[ht]
\centering
\caption{SFLA fault location accuracy across 80 scenarios.}
\label{tab:sfla}
\begin{tabular}{lcc}
\toprule
Metric & Normal comm.\ & Degraded comm.\ \\
\midrule
Mean location error (\% of line) & 0.8\% & 2.1\% \\
Max location error (\% of line)  & 1.9\% & 4.8\% \\
Location within $\pm 5\%$        & 100\% & 96\% \\
\bottomrule
\end{tabular}
\end{table}

\section{Conclusion}
\label{sec:conclusion}

\begin{finding}[WAP on IBR-rich IEEE~39-bus: 97.5\% agreement]
The PMU-assisted WAB with SFLA achieves 78/80 = 97.5\% agreement on the
IEEE~39-bus with high IBR penetration, exceeding the $\geq 97\%$ criterion.
The SFLA algorithm maintains fault location accuracy $\leq 2.1\%$ even under
degraded communication with 80\,ms latency.
\end{finding}

\begin{finding}[Islanded fallback is the weak point]
Under complete communication loss (islanded fallback), high-impedance faults
are not reliably detected by local Zone-2 elements alone. TR-77 (PMU state
estimation) addresses this gap via a stored state observer that can maintain
WAP function for up to 5\,s without live PMU data.
\end{finding}

\begin{thebibliography}{99}
\bibitem{Salim2008}
R.~H.~Salim \emph{et al.},
``Further developments on complex-valued neural network based fault location
for transmission lines with distributed generation,''
\emph{IET Gener.\ Transm.\ Distrib.}, vol.~5, no.~4, pp.~484--491, 2011.

\bibitem{Phadke2008}
A.~G.~Phadke and J.~S.~Thorp,
\emph{Synchronized Phasor Measurements and Their Applications},
Springer, New York, 2008.

\bibitem{IEC61968}
IEC~61968-1:2012, \emph{Application Integration at Electric Utilities ---
System Interfaces for Distribution Management}, IEC, Geneva, 2012.

\bibitem{IEEE1588}
IEEE~1588-2008, \emph{IEEE Standard for a Precision Clock Synchronization
Protocol for Networked Measurement and Control Systems}, IEEE, 2008.

\bibitem{Horowitz2008}
S.~H.~Horowitz and A.~G.~Phadke,
\emph{Power System Relaying}, 3rd ed., Wiley, Chichester, 2008.
\end{thebibliography}

\end{document}
